\section{Efficiency on a Low-Power Laptop}
\label{sec:efficiency}

All timings in this section were measured on the deployment target, a laptop with an Intel Core
i3-1115G4 (\V{hw.cores} cores, \V{hw.threads} threads), integrated graphics and \V{hw.ramgib}\,GiB
of memory, in Chrome \V{b.chrome} and in Python onnxruntime.

\paragraph{Quantization.}
Table~\ref{tab:quant} (Appendix~\ref{app:quant}) applies post-training quantization to the v1 models
\citep{jacob2018quantization,nagel2021whitepaper}. fp16 halves every file at no measurable cost
in accuracy, though on this CPU it is slightly slower than fp32 (hybrid: \V{q.hybrid.fp16.p50}
vs.\ \V{q.hybrid.onnxfp32.p50} ms). int8 always costs accuracy: dynamic quantization \V{q.cnn.int8dynamic.delta} to
\V{q.gru.int8dynamic.delta} points, static quantization \V{q.transformer.int8static.delta} to
\V{q.gru.int8static.delta} points, with the GRU hurt most. Its effect on latency depends on the
model. Dynamic int8 speeds up the transformer (p50 \V{q.transformer.int8dynamic.p50} vs.\
\V{q.transformer.fp32.p50} ms) but slows the CNN (\V{q.cnn.int8dynamic.p50} vs.\
\V{q.cnn.fp32.p50} ms) and the hybrid (\V{q.hybrid.int8dynamic.p50} vs.\
\V{q.hybrid.onnxfp32.p50} ms); both slowed models are built from depthwise convolutions, which
the transformer lacks, but we did not profile the cause. Every model is already fast enough in fp32
on this CPU, so the accuracy loss buys nothing; we ship the hybrid in fp16 (\V{m.ship.mb}\,MB).

\paragraph{Python vs.\ browser.}
One hybrid call takes \V{q.hybrid.fp16.p50} ms (p50) in Python onnxruntime with one thread and
\V{b.wasm1.p50} ms (p95 \V{b.wasm1.p95}) in onnxruntime-web with WebAssembly SIMD and one thread.
More threads do not help in the browser (two threads: p50 \V{b.wasm2.p50} ms; four:
\V{b.wasm4.p50} ms), and the WebGPU backend on the integrated GPU is slower still (p50
\V{b.webgpu1.p50} ms) while agreeing with WebAssembly on the top-1 class, so the app uses one
WebAssembly thread and keeps WebGPU as an option it adopts only if it measures faster. These are
the shipped file's numbers; Table~\ref{tab:models} and Figure~\ref{fig:acclat} come from the earlier
selection benchmark of every candidate (fp32, p50 \V{m.d.p50} ms for the hybrid), and the
per-segment times below are measured inside the running app, so the three differ by a few
milliseconds.

\paragraph{Live throughput.}
In the full app, MediaPipe runs in one worker on the GPU delegate and the classifier in another, so
a classification never blocks tracking. With the recorded clip as camera, tracking ran at
\V{t.hudon.fps} frames per second with the hand skeleton drawn (two \V{t.duration}-second runs:
\V{t.hudon.fps1} and \V{t.hudon.fps2}) and \V{t.hudoff.fps} without it (\V{t.hudoff.fps1},
\V{t.hudoff.fps2}); the order of the runs moved fps more than the overlay did. The classifier took
\V{t.hudon.cls1} and \V{t.hudon.cls2} ms per segment (p50), well inside the \V{budget.ms} ms
budget, and a word appeared \V{t.hudon.e2e1} ms (p50) after the segmenter decided that the sign
had ended; that decision itself comes \V{seg.lostms}--\V{seg.stillms} ms after the hands leave or
stop, so a user waits about \V{t.userlo}--\V{t.userhi} ms. Without a working GPU delegate the hand
landmarker alone takes \V{t.cpu.hand} ms per frame on the CPU and tracking drops to \V{t.cpu.fps}
fps, below our \V{budget.fps} fps target and below the frame rate of our own recordings.
